\chapterimage{figs/chapterhead.png}
\chapter{User Installation}
\label{chap:user-installation}

\section{Introduction}
\label{sec:user-installation-introduction}

This chapter explains how to install and start GenESyS as a user-facing
application, without requiring knowledge of the source tree or the internal
build layout. The focus is on supported installation paths, the runtime files
that appear after installation, the first launch of the GUI, and the checks
that confirm the program is ready for use.

The documented baseline for the repository is Ubuntu 24.04. The current
packaging and build surface are Debian-based and Qt6-based, so this chapter is
written around that environment. Other systems may work if they satisfy the
same dependencies, but they are not the baseline validated by this repository.

For a user, the important distinction is simple:
\begin{itemize}
\item installation puts a runnable GenESyS application on the system;
\item development build instructions are for contributors and validators.
\end{itemize}

The chapter therefore keeps the user path separate from the source-build path,
even though both can produce the same GUI executable name.

The visual sequence for this chapter is anchored by Figure~\ref{fig:user-installation-flow}
and then continues through Figure~\ref{fig:user-installed-structure},
Figure~\ref{fig:user-first-start},
Figure~\ref{fig:user-version-confirmation},
Figure~\ref{fig:user-executable-locations}, and
Figure~\ref{fig:user-installation-structure-summary}.

\begin{figure}[htbp]
    \centering

    % FIGURE-SPEC-BEGIN
    % Figure ID: user-installation-flow
    % Type: process diagram
    % Purpose:
    %   Show the supported installation and first-start sequence for users.
    % Source of truth:
    %   - debian/control
    %   - debian/rules
    %   - debian/genesys-simulator.install
    %   - source/applications/launcher/wrappers/launcher-wrapper.in
    %   - CMakePresets.json
    %   - packaging/linux/io.github.rlcancian.genesys.desktop
    % Required visual content:
    %   - supported Debian package path (single genesys-simulator package)
    %   - optional local artifact path
    %   - installed stable genesys-gui command name
    %   - launcher selecting per-user runtime or system fallback
    %   - first GUI startup
    %   - About-dialog version check
    % Accessibility description:
    %   Flowchart from package acquisition to a successful GUI launch.
    % Validation:
    %   Every node must match the current packaging and preset files.
    % FIGURE-SPEC-END

    \figureplaceholder{figs/generated/user-installation-flow}

    \caption{Supported installation and first-start flow for GenESyS.}
    \label{fig:user-installation-flow}
\end{figure}

\section{Supported installation paths}
\label{sec:user-installation-supported-paths}

The repository currently documents two practical ways to get a runnable
GenESyS GUI on the machine.

\subsection{Preferred path: Debian package}
\label{subsec:user-installation-debian-package}

GenESyS is distributed as a single Debian binary package,
\texttt{genesys-simulator}. Installing it is sufficient to obtain the complete
supported base: the Qt GUI (\texttt{genesys-gui}), the command-line shell
(\texttt{genesys-shell}), the HTTP worker (\texttt{genesys-worker}, with its
legacy \texttt{genesys-web} command alias), the stable launcher and dispatcher
mechanism, and the administrative update policy. GUI, Shell, and Worker remain
separate executable applications; they are simply not separate Debian
packages, so \texttt{apt install genesys-simulator} is enough to obtain a
working GUI, shell, and worker installation.

The Debian metadata installs a small, stable command at
\path{/usr/bin/genesys-gui} and also ships the desktop integration files under
the standard XDG locations:
\begin{itemize}
\item \path{/usr/share/applications/io.github.rlcancian.genesys.desktop};
\item \path{/usr/share/metainfo/io.github.rlcancian.genesys.metainfo.xml};
\item \path{/usr/share/icons/hicolor/scalable/apps/io.github.rlcancian.genesys.svg}.
\end{itemize}

\path{/usr/bin/genesys-gui} does not contain the GUI implementation itself.
It is a wrapper that starts the stable launcher mechanism installed by the
\texttt{genesys-simulator} package, which then starts the actual GUI: either a validated
runtime the user installed in their own home directory, or the
Debian-provided system installation if no per-user runtime is active. Section
\ref{sec:user-installation-launcher-and-runtime-selection} explains this in
detail; from a user's point of view, the important fact is that the same
\texttt{genesys-gui} command keeps working across that choice.

This is the clearest end-user route when a package artifact is available. It
creates the expected desktop launcher and provides a direct command-line entry
point.

If the package is available as a local file, install it with \texttt{apt} so
dependency resolution happens automatically:

\begin{lstlisting}[language=bash,caption={Installing the local GenESyS Debian package},label={lst:user-installation-local-deb}]
sudo apt install ./genesys-simulator_*.deb
\end{lstlisting}

If the package is already published in a configured repository, the same binary
package can be installed by name:

\begin{lstlisting}[language=bash,caption={Installing GenESyS from a repository},label={lst:user-installation-repo-deb}]
sudo apt install genesys-simulator
\end{lstlisting}

\subsection{Manual artifact installation}
\label{subsec:user-installation-manual-artifact}

When a binary artifact is delivered outside an apt repository, the same package
format should still be preferred. In practice, that means a user or maintainer
can download the \texttt{.deb} file and install it locally with the command
above.

If the artifact is provided as a staged filesystem tree rather than a package,
the repository does not define a separate user-facing installer for it. In that
case, the artifact should be treated as a validation or distribution artifact
that must already match the expected layout before it is copied into place.
This chapter does not define a new custom installer.

\section{System baseline and runtime dependencies}
\label{sec:user-installation-runtime-dependencies}

The current supported baseline is the Qt6 desktop stack on Ubuntu 24.04.
The Debian package metadata confirms the build-time requirements used by the
project:
\begin{itemize}
\item \texttt{cmake} 3.24 or newer;
\item \texttt{ninja-build};
\item \texttt{g++};
\item \texttt{pkgconf};
\item \texttt{qt6-base-dev};
\item \texttt{qt6-base-dev-tools};
\item \texttt{libgl-dev}.
\end{itemize}

For a packaged install, the user does not usually install those build-time
packages directly. They matter here because they define the toolchain that
produces the user package and the GUI runtime.

At runtime, the GUI expects a normal graphical desktop session, access to the
display server, and the standard desktop integration files provided by the
package. The application is intended to launch as a desktop GUI, not as a
headless background service.

\begin{figure}[htbp]
    \centering

    % FIGURE-SPEC-BEGIN
    % Figure ID: user-installed-structure
    % Type: filesystem layout diagram
    % Purpose:
    %   Show what is installed by the desktop package and where the launcher
    %   and metadata files live.
    % Source of truth:
    %   - debian/genesys-simulator.install
    %   - packaging/linux/io.github.rlcancian.genesys.desktop
    %   - packaging/linux/io.github.rlcancian.genesys.metainfo.xml
    % Required visual content:
    %   - /usr/bin/genesys-gui public wrapper
    %   - /usr/libexec/genesys/genesys-launcher and system fallback binary
    %   - desktop entry file
    %   - metainfo file
    %   - icon file
    % Accessibility description:
    %   Tree of installed files and desktop integration assets.
    % Validation:
    %   Paths must match the current package metadata.
    % FIGURE-SPEC-END

    \figureplaceholder{figs/generated/user-installed-structure}

    \caption{Installed GenESyS desktop files and launcher locations.}
    \label{fig:user-installed-structure}
\end{figure}

\section{Launcher, system fallback, and per-user runtime}
\label{sec:user-installation-launcher-and-runtime-selection}

The Debian package installs a stable base that works completely on its own,
with no per-user setup required. \path{/usr/bin/genesys-gui},
\path{/usr/bin/genesys-shell}, \path{/usr/bin/genesys-worker}, the legacy
\path{/usr/bin/genesys-web} alias, and \path{/usr/bin/genesys-mcp} are all
small wrapper commands. \texttt{genesys-gui} starts the stable launcher;
the other four start the stable, non-interactive dispatcher. Both mechanisms
live under \path{/usr/libexec/genesys/} and are installed by the
\texttt{genesys-simulator} package.

At each invocation, the launcher or dispatcher decides which runtime actually
serves the request:
\begin{itemize}
\item a validated runtime installed in the user's own home directory, under
  \path{~/.local/share/genesys/apps/<version>/}, selected through the
  \path{~/.local/share/genesys/current} symbolic link, if one is present and
  passes validation; or otherwise
\item the Debian-provided \emph{system fallback} runtime, installed under
  \path{/usr/libexec/genesys/system/bin/} by the \texttt{genesys-simulator}
  package.
\end{itemize}

A per-user runtime can therefore be updated without administrator privileges,
without ever touching a file owned by the Debian package, and it can always
fall back to the system-provided version if it is missing, incomplete, or
disabled by policy. Only \texttt{genesys-gui} is allowed to check for and
install an updated per-user runtime; the shell, worker, legacy web alias, and
MCP commands only dispatch to whichever runtime is already selected and never
perform network activity themselves.

\subsection{Administrative and per-user configuration}
\label{subsec:user-installation-update-configuration}

Two configuration files control this behavior:
\begin{itemize}
\item \path{/etc/genesys/update.conf} -- the system-wide administrative
  policy installed by the \texttt{genesys-simulator} package as a Debian
  conffile. It takes precedence over the per-user configuration.
\item \path{~/.config/genesys/update.conf} -- an optional per-user override,
  read only where the administrative policy allows it.
\end{itemize}

As shipped, the administrative policy disables remote update checks and
requires signature verification for any future update, so no network activity
happens unless a maintainer deliberately provisions and enables it. This is a
deliberate fail-closed default: the base application is fully usable from the
Debian package alone, and enabling remote per-user updates is a separate,
explicit administrative decision, not a side effect of installing the
package.

\subsection{Runtime status and troubleshooting}
\label{subsec:user-installation-launcher-log}

The launcher and dispatcher record their activity in
\path{~/.cache/genesys/launcher.log}. If \texttt{genesys-gui} unexpectedly
starts the system-provided runtime instead of an installed per-user runtime
(or vice-versa), this log is the first place to check; the current
\path{~/.local/share/genesys/current} link and the administrative
\path{allow_user_runtime} setting are the two most common causes of an
unexpected selection.

\section{First launch of the GUI}
\label{sec:user-installation-first-launch}

After installation, start the GUI from the desktop launcher or from a shell.
The installed binary name is \texttt{genesys-gui}. A terminal launch looks like
this:

\begin{lstlisting}[language=bash,caption={Starting the installed GenESyS GUI},label={lst:user-installation-start-gui}]
genesys-gui
\end{lstlisting}

If you are validating a local build rather than a package install, the current
source tree can produce the same executable through the \texttt{gui-app}
preset:

\begin{lstlisting}[language=bash,caption={Local build-and-run validation path},label={lst:user-installation-build-path}]
cmake --preset gui-app
cmake --build --preset gui-app --parallel "$(nproc)"
./build/gui-app/source/applications/gui/genesys/genesys-gui
\end{lstlisting}

The local build path is useful for validation, but it is not the same thing as
the user-install path. In the book, the package install is the primary user
path.

\begin{figure}[htbp]
    \centering

    % FIGURE-SPEC-BEGIN
    % Figure ID: user-first-start
    % Type: UI screenshot placeholder
    % Purpose:
    %   Show the GUI immediately after first launch.
    % Source of truth:
    %   - source/applications/gui/genesys/mainwindow.ui
    %   - source/applications/gui/genesys/mainwindow.cpp
    %   - source/applications/gui/genesys/controllers/DialogUtilityController.cpp
    % Required visual content:
    %   - main window
    %   - menu bar
    %   - workspace panels
    %   - About menu entry
    % Accessibility description:
    %   Main GenESyS window at startup.
    % Validation:
    %   The window must match the current GUI composition.
    % FIGURE-SPEC-END

    \figureplaceholder{figs/generated/user-first-start}

    \caption{First successful start of the GenESyS GUI.}
    \label{fig:user-first-start}
\end{figure}

\subsection{What to confirm on the first run}
\label{subsec:user-installation-first-run-checks}

The first launch should confirm only the basics:
\begin{itemize}
\item the window opens without an immediate failure;
\item the menus and toolbars are visible;
\item the main workspace loads normally;
\item the application can be closed cleanly.
\end{itemize}

The user does not need to build a model yet. The goal is only to prove that the
installed application starts and displays the main GUI.

\subsection{Version confirmation}
\label{subsec:user-installation-version-confirmation}

The GUI includes an \textit{About} menu with an \textit{About Genesys} action.
That action is the most direct way to confirm the application identity and
version information from inside the running program.

\begin{figure}[htbp]
    \centering

    % FIGURE-SPEC-BEGIN
    % Figure ID: user-version-confirmation
    % Type: dialog screenshot placeholder
    % Purpose:
    %   Show the About dialog used to confirm the runtime version.
    % Source of truth:
    %   - source/applications/gui/genesys/mainwindow.ui
    %   - source/applications/gui/genesys/controllers/DialogUtilityController.cpp
    % Required visual content:
    %   - About Genesys dialog
    %   - version field or equivalent runtime identity
    %   - application name
    % Accessibility description:
    %   About dialog used for version and identity confirmation.
    % Validation:
    %   The dialog content must match the current GUI code.
    % FIGURE-SPEC-END

    \figureplaceholder{figs/generated/user-version-confirmation}

    \caption{Version confirmation through the GenESyS About dialog.}
    \label{fig:user-version-confirmation}
\end{figure}

\section{Where the files live after installation}
\label{sec:user-installation-file-locations}

From a user point of view, the important installed locations are the stable
public commands and the desktop integration files. For package-based
installs, the key paths are:
\begin{itemize}
\item \path{/usr/bin/genesys-gui}, \path{/usr/bin/genesys-shell},
  \path{/usr/bin/genesys-worker} -- stable public commands, never removed or
  replaced by a per-user update;
\item \path{/usr/libexec/genesys/} -- the launcher, the dispatcher, and the
  Debian-provided system fallback GUI/Shell/Worker executables, under
  \texttt{system/bin/};
\item \path{/etc/genesys/update.conf} -- administrative configuration;
\item \path{/usr/share/applications/io.github.rlcancian.genesys.desktop};
\item \path{/usr/share/metainfo/io.github.rlcancian.genesys.metainfo.xml};
\item \path{/usr/share/icons/hicolor/scalable/apps/io.github.rlcancian.genesys.svg}.
\end{itemize}

If a per-user runtime has been installed, it lives entirely under the user's
own home directory and is never mixed with the paths above:
\begin{itemize}
\item \path{~/.local/share/genesys/apps/<version>/} -- one directory per
  installed runtime version;
\item \path{~/.local/share/genesys/current} -- symbolic link to the active
  version;
\item \path{~/.config/genesys/update.conf} -- optional per-user
  configuration;
\item \path{~/.cache/genesys/launcher.log} -- launcher/dispatcher activity
  log.
\end{itemize}

For local validation from source, the matching build-tree executable is:
\begin{itemize}
\item \path{build/gui-app/source/applications/gui/genesys/genesys-gui}.
\end{itemize}

\begin{figure}[htbp]
    \centering

    % FIGURE-SPEC-BEGIN
    % Figure ID: user-executable-locations
    % Type: layout diagram
    % Purpose:
    %   Show the installed binary and the equivalent local build-tree path.
    % Source of truth:
    %   - debian/genesys-simulator.install
    %   - CMakePresets.json
    %   - packaging/linux/io.github.rlcancian.genesys.desktop
    % Required visual content:
    %   - installed binary path
    %   - desktop launcher path
    %   - build-tree executable path
    % Accessibility description:
    %   Side-by-side installed and build-tree executable locations.
    % Validation:
    %   Paths must be current and directly runnable where applicable.
    % FIGURE-SPEC-END

    \figureplaceholder{figs/generated/user-executable-locations}

    \caption{Installed and build-tree executable locations for GenESyS.}
    \label{fig:user-executable-locations}
\end{figure}

\section{Updating the installation}
\label{sec:user-installation-updating}

If GenESyS was installed from a Debian package, update it by installing the new
package version over the old one. On a configured package repository, the
standard flow is:

\begin{lstlisting}[language=bash,caption={Updating an apt-based GenESyS installation},label={lst:user-installation-update}]
sudo apt update
sudo apt install --only-upgrade genesys-simulator
\end{lstlisting}

If you are using a local \texttt{.deb} file, replace the package from the new
artifact:

\begin{lstlisting}[language=bash,caption={Replacing a local package artifact},label={lst:user-installation-update-local}]
sudo apt install ./genesys-simulator_*.deb
\end{lstlisting}

This keeps the package manager in control of file ownership and dependency
tracking. It is preferable to copying files by hand into system directories.

\section{Removing the installation}
\label{sec:user-installation-removal}

To remove the package, use the package manager instead of deleting files
manually. A plain remove keeps configuration if the package has any:

\begin{lstlisting}[language=bash,caption={Removing the GenESyS package},label={lst:user-installation-remove}]
sudo apt remove genesys-simulator
\end{lstlisting}

If a full purge is desired, remove configuration files as well:

\begin{lstlisting}[language=bash,caption={Purging the GenESyS package},label={lst:user-installation-purge}]
sudo apt purge genesys-simulator
\end{lstlisting}

A single \texttt{genesys-simulator} package owns every installed file, so one
remove or purge command is enough to uninstall GenESyS completely; there is no
separate package name to track down afterward. The historical
\path{genesys-web} command remains available as a compatibility alias for
\path{genesys-worker} inside the same package. The historical development
package names (\texttt{genesys-common}, \texttt{genesys-shell},
\texttt{genesys-worker}, \texttt{genesys-gui}, and \texttt{genesys-web}) are
declared as \texttt{Breaks}/\texttt{Replaces} only so that a machine with an
earlier development installation can upgrade cleanly to
\texttt{genesys-simulator}; they are not produced as binary packages anymore.

Removing or purging the Debian package never touches a per-user runtime
installed under \path{~/.local/share/genesys/}. That directory, and anything
in it, is owned by the user, not by \texttt{dpkg}, and remains after the
package is gone. Deleting it, if ever desired, is a separate, manual step.

\section{Initial troubleshooting}
\label{sec:user-installation-troubleshooting}

If the GUI does not start, the first checks should be practical rather than
architectural:
\begin{itemize}
\item verify that the package installed successfully;
\item verify that \texttt{genesys-gui} is on the path or run from the full
  installed location;
\item verify that the desktop session can open Qt applications;
\item verify that the package dependencies were resolved;
\item verify that the display session is available if the command is started
from a terminal.
\end{itemize}

If the application starts but the About dialog or launcher does not look
correct, recheck the package version and the desktop integration files. If the
local build path is being used, confirm that the \texttt{gui-app} preset was
selected and that the build completed successfully before starting the binary.

\section{Limitations}
\label{sec:user-installation-limitations}

This chapter does not claim support for every operating system or every
distribution packaging style. It documents the Debian-based package path and
the local build-tree validation path that are currently confirmed in the
repository.

It also does not redefine the installer format. If a future release adds a
different distribution artifact, that artifact should be documented with the
same discipline as the Debian package: package name, installed paths, launch
command, update path, removal path, and first-run verification.

\begin{figure}[htbp]
    \centering

    % FIGURE-SPEC-BEGIN
    % Figure ID: user-installation-structure-summary
    % Type: summary diagram
    % Purpose:
    %   Summarize the relationship between install, launch, verification,
    %   update, and removal.
    % Source of truth:
    %   - debian/control
    %   - debian/genesys-gui.install
    %   - CMakePresets.json
    %   - packaging/linux/io.github.rlcancian.genesys.desktop
    % Required visual content:
    %   - install
    %   - launch
    %   - About check
    %   - update
    %   - remove
    % Accessibility description:
    %   Summary lifecycle for the user package.
    % Validation:
    %   All nodes must correspond to commands and files already documented.
    % FIGURE-SPEC-END

    \figureplaceholder{figs/generated/user-installation-structure-summary}

    \caption{Lifecycle summary for a user installation of GenESyS.}
    \label{fig:user-installation-structure-summary}
\end{figure}

\section{Final considerations}
\label{sec:user-installation-final-considerations}

The user installation path for GenESyS is intentionally narrow: install the
desktop package when available, or use the local build-tree validation path
only when you are checking the runtime from source. In both cases, the outcome
that matters is the same: the \texttt{genesys-gui} application starts, displays
its main window, and exposes the expected About dialog and desktop launcher.

Once that works, the reader can move to the next chapter and start learning how
the application is organized and what the main GUI areas mean in practice.

Test your understanding of this content by answering the following questions:

\begin{enumerate}
\item Why is the Debian package the preferred user-installation path in this chapter?
\item Which files and paths confirm that GenESyS was installed correctly on the system?
\item How does the chapter distinguish a user installation from a local build-tree validation?
\end{enumerate}
